\documentclass[a4paper,landscape,8pt]{extarticle}
\usepackage[landscape,margin=0.7cm,top=0.5cm,bottom=0.7cm]{geometry}
\usepackage[utf8]{inputenc}
\usepackage[T1]{fontenc}
\usepackage[ngerman]{babel}
\usepackage{helvet}
\renewcommand{\familydefault}{\sfdefault}
\usepackage{xcolor}
\usepackage{tikz}
\usetikzlibrary{positioning}
\usepackage{enumitem}
\usepackage{multicol}
\usepackage{array}
\usepackage{fancyhdr}
\usepackage{amssymb}

% Farben
\definecolor{simblue}{HTML}{1565C0}
\definecolor{devicegreen}{HTML}{2E7D32}
\definecolor{applered}{HTML}{C62828}
\definecolor{lightgray}{HTML}{F5F5F5}
\definecolor{bordergray}{HTML}{BBBBBB}

\pagestyle{fancy}
\fancyhf{}
\renewcommand{\headrulewidth}{0pt}
\fancyfoot[C]{\footnotesize Seite \thepage{} von 3}

\setlength{\parindent}{0pt}
\setlength{\parskip}{2pt}

% Kompakte Listen
\setlist{nosep, leftmargin=1em, itemsep=0pt, topsep=1pt}

\begin{document}

% ============================================================================
% SEITE 1: INFORMATIONSBLATT
% ============================================================================

\noindent{\Large\bfseries Die drei Schlüssel zu deinem iPhone} \hfill
{\footnotesize\textbf{Lernziele:} Codes unterscheiden | Schutzfunktionen | Folgen bei Verlust}

\vspace{1mm}
\hrule
\vspace{1mm}

% Drei-Spalten-Layout
\begin{minipage}[t]{0.32\linewidth}
\begin{center}
\colorbox{simblue!10}{%
\begin{minipage}{0.95\linewidth}
\vspace{2mm}
\begin{center}
{\Large\bfseries\textcolor{simblue}{1. SIM-PIN}}\\[1pt]
{\footnotesize\itshape Der Schlüssel zur Telefonnummer}
\end{center}
\vspace{2mm}
\end{minipage}%
}
\end{center}

\vspace{1mm}

\textbf{Was ist das?} Ein \textbf{4-stelliger Code} für die SIM-Karte in deinem Handy.

\textbf{Analogie:} Wie der \textbf{Schlüssel zum Briefkasten} -- gehört zur Telefonnummer, nicht zum Haus.

\textbf{Wann gebraucht?}
\begin{itemize}
\item Nach Neustart des Handys
\item Wenn SIM in anderes Gerät kommt
\end{itemize}

\textbf{Wer kennt ihn?} Dein \textbf{Mobilfunkanbieter} hat ihn gegeben. Standard: 0000 oder 1234.

\textbf{Fehlversuche:}
\begin{itemize}
\item 3× falsch → SIM gesperrt
\item Dann \textbf{PUK} nötig (8 Ziffern, auf Unterlagen)
\item 10× PUK falsch → SIM \textbf{kaputt}!
\end{itemize}

\textbf{Ohne ihn:} Keine Anrufe/SMS/mobiles Internet. \textit{Aber WLAN geht!}

\end{minipage}%
\hfill%
\begin{minipage}[t]{0.32\linewidth}
\begin{center}
\colorbox{devicegreen!10}{%
\begin{minipage}{0.95\linewidth}
\vspace{2mm}
\begin{center}
{\Large\bfseries\textcolor{devicegreen}{2. Geräte-Code}}\\[1pt]
{\footnotesize\itshape Der Schlüssel zum iPhone selbst}
\end{center}
\vspace{2mm}
\end{minipage}%
}
\end{center}

\vspace{1mm}

\textbf{Was ist das?} Ein \textbf{6-stelliger Code}, den \textbf{du selbst} gewählt hast.

\textbf{Analogie:} Wie der \textbf{Hausschlüssel} -- öffnet alles im iPhone.

\textbf{Face ID / Touch ID:} Nur \textbf{Abkürzungen}! Der Code bleibt der echte Schlüssel.

\textbf{Wann gebraucht?}
\begin{itemize}
\item Nach Neustart (Face ID erst danach)
\item Nach 48h ohne Nutzung
\item Für wichtige Einstellungen
\end{itemize}

\textbf{Was schützt er?}
\begin{itemize}
\item Alle Daten auf dem Gerät
\item \textbf{Schlüsselbund} (gespeicherte Passwörter)
\item Zugang zu Apps und Fotos
\end{itemize}

\textbf{Wer kennt ihn?} \textbf{Nur du!} Apple kennt ihn NICHT.

\textbf{Bei Verlust:} iPhone muss \textbf{komplett gelöscht} werden. Alle Daten weg (außer iCloud-Backup)!

\end{minipage}%
\hfill%
\begin{minipage}[t]{0.32\linewidth}
\begin{center}
\colorbox{applered!10}{%
\begin{minipage}{0.95\linewidth}
\vspace{2mm}
\begin{center}
{\Large\bfseries\textcolor{applered}{3. Apple-ID}}\\[1pt]
{\footnotesize\itshape Der Generalschlüssel bei Apple}
\end{center}
\vspace{2mm}
\end{minipage}%
}
\end{center}

\vspace{1mm}

\textbf{Was ist das?} Dein \textbf{Konto bei Apple} -- sieht aus wie E-Mail, ist aber mehr!

\textbf{Analogie:} Wie \textbf{Ausweis + Bankkarte} bei einer Firma, die alles aufbewahrt.

\textbf{Was ist verbunden?}
\begin{itemize}
\item iCloud (Fotos, Kontakte, Backup)
\item App Store (gekaufte Apps)
\item \textbf{„Wo ist?"} -- iPhone orten
\item iMessage, FaceTime
\end{itemize}

\textbf{Mächtigste Funktion:} Von \textit{jedem} Computer iPhone \textbf{orten, sperren, löschen}!

\textbf{Wer kennt es?} \textbf{Nur du!} Apple kann zurücksetzen, aber nicht sehen.

\textbf{Zwei-Faktor:} Bei Anmeldung kommt Code aufs iPhone. Ohne iPhone = kein Zugang.

\textbf{WARNUNG:} Wer Apple-ID + Passwort hat, kann \textbf{alles sehen und löschen}!

\end{minipage}

\vspace{1mm}

% Zusammenfassung
\begin{center}
\fcolorbox{bordergray}{lightgray}{%
\begin{minipage}{0.98\linewidth}
\centering
\vspace{1.5mm}
\textbf{Zusammenfassung:}
\textcolor{simblue}{\textbf{SIM-PIN}} = Telefonnummer (vom Anbieter) ~$|$~
\textcolor{devicegreen}{\textbf{Geräte-Code}} = iPhone-Inhalt (nur du) ~$|$~
\textcolor{applered}{\textbf{Apple-ID}} = Alles bei Apple + Fernzugriff
\vspace{1.5mm}
\end{minipage}%
}
\end{center}

\newpage

% ============================================================================
% SEITE 2: ÜBUNGEN
% ============================================================================

\begin{center}
{\LARGE\bfseries Übungen: Teste dein Wissen}\\[3pt]
{\normalsize Schau dir Seite 1 an und beantworte die Fragen}
\end{center}

\vspace{3mm}
\hrule
\vspace{4mm}

\begin{multicols}{2}

\textbf{\large Teil A: Zuordnung}

Welcher Schlüssel ist gemeint? Schreibe: \textit{SIM-PIN}, \textit{Geräte-Code} oder \textit{Apple-ID}

\vspace{2mm}

\renewcommand{\arraystretch}{1.8}
\begin{tabular}{@{}p{10.5cm}p{3.5cm}@{}}
1. Hat 4 Ziffern und kommt vom Mobilfunkanbieter: & \dotfill \\
2. Hat 6 Ziffern und hast du selbst gewählt: & \dotfill \\
3. Sieht aus wie eine E-Mail-Adresse: & \dotfill \\
4. Face ID ist eine Abkürzung dafür: & \dotfill \\
5. Damit kannst du dein iPhone aus der Ferne löschen: & \dotfill \\
6. Wenn er 3× falsch eingegeben wird, brauchst du den PUK: & \dotfill \\
\end{tabular}
\renewcommand{\arraystretch}{1}

\vspace{6mm}

\textbf{\large Teil B: Richtig oder Falsch?}

Kreuze an: R = richtig, F = falsch

\vspace{2mm}

\renewcommand{\arraystretch}{1.6}
\begin{tabular}{@{}p{10cm}cc@{}}
& R & F \\
1. Apple kennt meinen Geräte-Code. & $\square$ & $\square$ \\
2. Die SIM-PIN schützt meine Fotos. & $\square$ & $\square$ \\
3. Ohne SIM-PIN kann ich noch WLAN nutzen. & $\square$ & $\square$ \\
4. Die Apple-ID brauche ich, um Apps zu kaufen. & $\square$ & $\square$ \\
5. Touch ID ersetzt den Geräte-Code komplett. & $\square$ & $\square$ \\
6. Mit der Apple-ID kann ich mein verlorenes iPhone finden. & $\square$ & $\square$ \\
\end{tabular}
\renewcommand{\arraystretch}{1}

\columnbreak

\textbf{\large Teil C: Situationen}

Was tust du? Schreibe die Antwort in Stichworten.

\vspace{3mm}

\textbf{Situation 1:}\\
Dein iPhone zeigt „SIM gesperrt -- PUK eingeben". Du weißt den PUK nicht.

\textit{Was machst du?}
\vspace{8mm}

\rule{\linewidth}{0.4pt}
\vspace{5mm}

\rule{\linewidth}{0.4pt}

\vspace{5mm}

\textbf{Situation 2:}\\
Du hast deinen Geräte-Code vergessen. Face ID geht auch nicht.

\textit{Was passiert mit deinen Daten auf dem iPhone?}
\vspace{8mm}

\rule{\linewidth}{0.4pt}
\vspace{5mm}

\rule{\linewidth}{0.4pt}

\vspace{5mm}

\textbf{Situation 3:}\\
Jemand fragt nach deinem Apple-ID-Passwort „um dir zu helfen".

\textit{Was solltest du tun?}
\vspace{8mm}

\rule{\linewidth}{0.4pt}
\vspace{5mm}

\rule{\linewidth}{0.4pt}

\vspace{8mm}

\textbf{\large Teil D: Die drei Ebenen}

Male drei Kreise ineinander (wie eine Zielscheibe) und beschrifte sie:

\begin{itemize}
\item \textbf{Innen} (kleinster): Was nur die Telefon-Funktion betrifft
\item \textbf{Mitte}: Was auf dem Gerät gespeichert ist
\item \textbf{Außen} (größter): Was auch im Internet bei Apple liegt
\end{itemize}

\end{multicols}

\newpage

% ============================================================================
% SEITE 3: CHECKLISTE
% ============================================================================

\begin{center}
{\LARGE\bfseries Checkliste: Das habe ich verstanden}\\[3pt]
{\normalsize Geh die Liste durch und hake ab, was du sicher weißt}
\end{center}

\vspace{3mm}
\hrule
\vspace{6mm}

% Drei Boxen nebeneinander
\begin{minipage}[t]{0.31\linewidth}
\begin{center}
\fcolorbox{simblue}{simblue!8}{%
\begin{minipage}{0.92\linewidth}
\vspace{3mm}
\begin{center}
{\large\bfseries\textcolor{simblue}{SIM-PIN}}
\end{center}
\vspace{2mm}
\begin{itemize}[label=$\square$, itemsep=5pt]
\item Ich weiß, dass die SIM-PIN vom \textbf{Mobilfunkanbieter} kommt
\item Ich verstehe, dass sie nur die \textbf{Telefonnummer} schützt
\item Ich weiß, wo ich meinen \textbf{PUK} finden kann
\item Ich weiß: 3× falsch = SIM gesperrt
\item Ich weiß: 10× PUK falsch = SIM kaputt
\end{itemize}
\vspace{3mm}
\end{minipage}%
}
\end{center}
\end{minipage}%
\hfill%
\begin{minipage}[t]{0.31\linewidth}
\begin{center}
\fcolorbox{devicegreen}{devicegreen!8}{%
\begin{minipage}{0.92\linewidth}
\vspace{3mm}
\begin{center}
{\large\bfseries\textcolor{devicegreen}{Geräte-Code}}
\end{center}
\vspace{2mm}
\begin{itemize}[label=$\square$, itemsep=5pt]
\item Ich weiß, dass \textbf{ich} diesen Code selbst gewählt habe
\item Ich verstehe: Face ID ist nur eine \textbf{Abkürzung}
\item Ich weiß, dass Apple diesen Code \textbf{nicht} kennt
\item Ich verstehe: Vergessen = iPhone \textbf{löschen}
\item Ich weiß, was der \textbf{Schlüsselbund} ist
\end{itemize}
\vspace{3mm}
\end{minipage}%
}
\end{center}
\end{minipage}%
\hfill%
\begin{minipage}[t]{0.31\linewidth}
\begin{center}
\fcolorbox{applered}{applered!8}{%
\begin{minipage}{0.92\linewidth}
\vspace{3mm}
\begin{center}
{\large\bfseries\textcolor{applered}{Apple-ID}}
\end{center}
\vspace{2mm}
\begin{itemize}[label=$\square$, itemsep=5pt]
\item Ich verstehe: Apple-ID = mein \textbf{Konto} bei Apple
\item Ich weiß, was alles damit verbunden ist
\item Ich kann mein iPhone damit \textbf{orten und löschen}
\item Ich verstehe den \textbf{Zwei-Faktor-Schutz}
\item Ich gebe das Passwort \textbf{niemals} weiter
\end{itemize}
\vspace{3mm}
\end{minipage}%
}
\end{center}
\end{minipage}

\vspace{8mm}

% Gesamtverständnis
\begin{center}
\fcolorbox{black}{lightgray}{%
\begin{minipage}{0.95\linewidth}
\vspace{4mm}
\begin{center}
{\large\bfseries Gesamtverständnis}
\end{center}
\vspace{3mm}
\begin{multicols}{2}
\begin{itemize}[label=$\square$, itemsep=6pt]
\item Ich kann die \textbf{drei Codes unterscheiden} und weiß, welcher wofür ist
\item Ich verstehe, dass \textbf{SIM-PIN} und \textbf{Geräte-Code} verschiedene Dinge schützen
\item Ich weiß, dass die \textbf{Apple-ID} die mächtigste ist (Fernzugriff!)
\item Ich verstehe, warum ich meine \textbf{Passwörter sicher aufbewahren} sollte
\item Ich weiß, an wen ich mich wende, wenn ich einen Code vergesse
\item Ich fühle mich \textbf{sicherer} im Umgang mit meinem iPhone
\end{itemize}
\end{multicols}
\vspace{3mm}
\end{minipage}%
}
\end{center}

\vspace{8mm}

% Notfall-Notizen
\begin{center}
\fcolorbox{bordergray}{white}{%
\begin{minipage}{0.95\linewidth}
\vspace{4mm}
\begin{center}
{\large\bfseries Meine Notfall-Notizen}\\[3pt]
{\footnotesize (Fülle das hier aus und bewahre es sicher auf!)}
\end{center}
\vspace{4mm}
\renewcommand{\arraystretch}{2}
\begin{tabular}{@{}p{0.3\linewidth}p{0.3\linewidth}p{0.3\linewidth}@{}}
Mein Mobilfunkanbieter: & Wo ist mein PUK? & Wer hilft mir bei Problemen? \\
\dotfill & \dotfill & \dotfill \\[3mm]
Meine Apple-ID (E-Mail): & & \\
\dotfill & & \\
\end{tabular}
\renewcommand{\arraystretch}{1}
\vspace{4mm}
\end{minipage}%
}
\end{center}

\vfill

\begin{center}
\footnotesize\textit{Stand: \today{} -- Bei Fragen: Frag jemanden, dem du vertraust! Gib niemals Passwörter an Fremde weiter.}
\end{center}

\end{document}
